% MATH 147 -- Zorich-style rewrite (experiment).
% Same qhnotes package and options as ../note-course.
\documentclass[10pt]{article}
\usepackage[color]{qhnotes}

% Document-local shortcuts (same as ../note-course/note.tex)
\newcommand{\paren}[1]{\left(#1\right)}
\newcommand{\abs}[1]{\left|#1\right|}
\newcommand{\norm}[1]{\left\lVert#1\right\rVert}
\newcommand{\bracks}[1]{\left[#1\right]}
\newcommand{\set}[1]{\left\{#1\right\}}
\newcommand{\R}{\mathbb{R}}
\newcommand{\N}{\mathbb{N}}
\newcommand{\Z}{\mathbb{Z}}
\newcommand{\Q}{\mathbb{Q}}
\newcommand{\C}{\mathbb{C}}
\newcommand{\ds}{\displaystyle}

% Zorich-style headings: 1 Chapter / 1.1 Section / a. Subsection
\renewcommand{\thesubsubsection}{\alph{subsubsection}}
\makeatletter
\renewcommand{\@seccntformat}[1]{%
  \csname the#1\endcsname\csname qh@punct@#1\endcsname\quad}
\newcommand{\qh@punct@subsubsection}{.}
\makeatother

\title{\textbf{Intro to Analysis}\\[0.5em]\large MATH 147 \\[0.3em]
  \normalsize (written in the style of Zorich)}
\author{Qinghao Hu}
\date{\today}
\course{Intro to Analysis}

\begin{document}

\maketitle
\tableofcontents

%!TEX root = ../note.tex
\section{The Real Numbers}

Differential calculus studies slopes and tangent lines; integral calculus
studies areas and accumulation. Both are built on the real numbers, and the
purpose of this chapter is to say exactly what the real numbers are and
which of their properties analysis actually uses.

%!TEX root = ../../note.tex
\subsection{Numbers and the Language of Logic}

\subsubsection{The number systems}
We start from the familiar systems
\begin{align*}
    \N &= \set{1,2,3,\ldots}, \\
    \Z &= \set{\ldots,-1,0,1,\ldots}, \\
    \Q &= \set{\frac{m}{n} : m,n\in\Z,\ n\ne 0}.
\end{align*}
Each one repairs a defect of the previous one. $\N$ is closed under
addition; $\Z$ is also closed under subtraction and multiplication, but not
under division ($1/2\notin\Z$); $\Q$ admits all quotients with nonzero
denominator. Even $\Q$, however, is too small: it has no number whose square
is $2$.

\begin{theorem}[Irrationality of $\sqrt2$]\label{thm:sqrt2-irrational}
    $\sqrt{2}\notin\Q$.
\end{theorem}

The idea is to write $\sqrt2$ in lowest terms; the equation $m^2=2n^2$ then
forces both $m$ and $n$ to be even.

\begin{proof}
    Suppose $\sqrt2=m/n$ with $m,n\in\Z$, $n>0$, $\gcd(m,n)=1$. Squaring
    gives $m^2=2n^2$, so $m^2$ is even. Since an odd integer has an odd
    square, $m$ is even, say $m=2\ell$. Then $4\ell^2=2n^2$, i.e.\
    $n^2=2\ell^2$, so $n$ is even by the same argument. Thus $2$ divides
    both $m$ and $n$, contradicting $\gcd(m,n)=1$.
\end{proof}

The elements of $\R\setminus\Q=\set{x\in\R : x\notin\Q}$ are called
\term{irrational numbers}.

\subsubsection{Decimal representation}
Informally, every nonnegative real number has a decimal expansion
\begin{equation*}
    x=a+\sum_{i=1}^{\infty}\frac{x_i}{10^i},
    \qquad a\in\N\cup\set{0},
    \qquad x_i\in\set{0,1,\ldots,9},
\end{equation*}
and negative numbers are negatives of such expansions. The expansion need
not be unique.

\begin{example}
    $1.999\cdots=2$. Indeed, if $x=1.999\cdots$, then $10x=19.999\cdots$, so
    $9x=10x-x=18$ and $x=2$.
\end{example}

This description of $\R$ is convenient but not a foundation: it does not
tell us what an infinite sum is. In \S\ref{sec:axioms} we replace it by a
list of axioms.

\subsubsection{Statements and quantifiers}
A \term{statement} is a sentence that is either true or false. We use the
\term{quantifiers} $\forall$ (``for all'') and $\exists$ (``there
exists''). Each quantifier is written in parentheses, and the statement it
governs in square brackets. For example,
\begin{align*}
    (\forall x\in\R)\,[x^2-x\geq 0] &\quad\text{is false (take } x=\tfrac12\text{)}, \\
    (\exists x\in\R)\,[x^2-x\geq 0] &\quad\text{is true (take } x=2\text{)}.
\end{align*}
Negation interchanges the quantifiers:
\begin{equation*}
    \neg\bigl[(\forall x)\,P(x)\bigr]\iff (\exists x)\,[\neg P(x)],
    \qquad
    \neg\bigl[(\exists x)\,P(x)\bigr]\iff (\forall x)\,[\neg P(x)].
\end{equation*}

The \emph{order} of quantifiers matters, as the next two propositions show.

\begin{proposition}\label{prop:forall-exists}
    $(\forall x\in\R)(\exists y\in\R)\,[x^3-y^3=1]$.
\end{proposition}

\begin{proof}
    Given $x$, take $y=\sqrt[3]{x^3-1}$; then $x^3-y^3=1$.
\end{proof}

\begin{proposition}
    The statement $(\exists y\in\R)(\forall x\in\R)\,[x^3-y^3=1]$ is false.
\end{proposition}

\begin{proof}
    We prove its negation $(\forall y\in\R)(\exists x\in\R)\,[x^3-y^3\neq1]$.
    Given $y$, take $x=y$; then $x^3-y^3=0\neq1$.
\end{proof}

\begin{remark}
    In Proposition~\ref{prop:forall-exists} the number $y$ may depend on
    $x$, and it does: $y=\sqrt[3]{x^3-1}$. The second statement demands a
    single $y$ that works for every $x$ at once, which is much stronger.
    Interchanging $\forall$ and $\exists$ can turn a true statement into a
    false one.
\end{remark}

%!TEX root = ../../note.tex
\subsection{The Natural Numbers: Induction and Well-Ordering}

Two properties of $\N$ are used throughout the course. Induction lets us
prove a statement for every $n$ by passing from $n$ to $n+1$; well-ordering
lets us pick the least element of a set of natural numbers. We take the first
as a principle and derive the second from it.

\subsubsection{Mathematical induction}
\begin{axiom}[Principle of Mathematical Induction]
    Let $P(n)$ be a statement depending on $n\in\N$. If
    \begin{enumerate}[label=(\arabic*)]
        \item $P(1)$ is true, and
        \item for every $k\in\N$, $P(k)$ implies $P(k+1)$,
    \end{enumerate}
    then $P(n)$ is true for every $n\in\N$.
\end{axiom}

Thus an induction proof has two parts: the \term{base case} $P(1)$, and the
\term{inductive step}, in which we fix an arbitrary $k$, assume $P(k)$ (the
\emph{inductive hypothesis}), and deduce $P(k+1)$.

\begin{example}
    $1+2+\cdots+n=\dfrac{n(n+1)}{2}$ for every $n\in\N$.
\end{example}

\begin{proof}
    For $n=1$ both sides equal $1$. If the formula holds for $k$, then
    \begin{equation*}
        1+\cdots+k+(k+1)
        = \frac{k(k+1)}{2}+(k+1)
        = \frac{k(k+1)+2(k+1)}{2}
        = \frac{(k+1)(k+2)}{2},
    \end{equation*}
    which is the formula for $k+1$.
\end{proof}

The same principle, stated for sets, reads as follows.

\begin{proposition}\label{prop:induction-sets}
    If $S\subseteq\N$, $1\in S$, and $k\in S\Rightarrow k+1\in S$, then
    $S=\N$.
\end{proposition}

\begin{proof}
    Apply induction to the statement $P(n)$: ``$n\in S$''.
\end{proof}

\subsubsection{The well-ordering principle}
\begin{theorem}[Well-Ordering Principle]\label{thm:wop}
    Every nonempty subset $S$ of $\N$ has a least element.
\end{theorem}

Suppose $S$ had no least element. We show by induction that then $S$
contains none of $1,2,\ldots,n$, for every $n$; hence $S=\emptyset$.

\begin{proof}
    Assume $S\neq\emptyset$ has no least element, and let
    \begin{equation*}
        T=\set{n\in\N : \set{1,2,\ldots,n}\cap S=\emptyset}.
    \end{equation*}
    \emph{$1\in T$}: otherwise $1\in S$, and $1$, being the least natural
    number, would be the least element of $S$.

    \emph{$k\in T\Rightarrow k+1\in T$}: if $\set{1,\ldots,k}\cap S=\emptyset$
    but $\set{1,\ldots,k+1}\cap S\neq\emptyset$, then $k+1\in S$ while no
    smaller number lies in $S$, so $k+1$ would be the least element of $S$.

    By Proposition~\ref{prop:induction-sets}, $T=\N$. Hence every $n\in\N$
    satisfies $n\notin S$, i.e.\ $S=\emptyset$, a contradiction.
\end{proof}

Well-ordering gives the method of the \term{minimal counterexample}: if
$P(n)$ failed for some $n$, the set $S=\set{n\in\N:\neg P(n)}$ would be
nonempty and would have a least element $k$. Then $P(k-1)$ holds by
minimality, and one shows that $P(k-1)$ forces $P(k)$ --- a contradiction.

\begin{example}
    Using well-ordering, prove that for every $n\in\N$
    \begin{equation}\label{eq:sum-squares}
        1^2 + 2^2 + \cdots + n^2 = \frac{n(n + 1)(2n + 1)}{6}.
    \end{equation}
\end{example}

\begin{proof}
    Let $P(n)$ be the statement \eqref{eq:sum-squares}, and suppose that
    the set of counterexamples
    \begin{equation*}
        S=\set{n \in \N : \neg P(n)}
    \end{equation*}
    is nonempty. By
    Theorem~\ref{thm:wop}, $S$ has a least element $k$. Since $P(1)$ holds,
    $k>1$, so $k-1\in\N$, and $P(k-1)$ holds because $k-1<k$:
    \begin{equation*}
        1^2 + \cdots + (k - 1)^2 = \frac{(k - 1)k(2k - 1)}{6}.
    \end{equation*}
    Adding $k^2$,
    \begin{align*}
        1^2 + \cdots + k^2
        &= \frac{k\bigl((k-1)(2k-1)+6k\bigr)}{6} \\
        &= \frac{k(2k^2+3k+1)}{6}
         = \frac{k(k+1)(2k+1)}{6},
    \end{align*}
    so $P(k)$ holds, contradicting $k\in S$. Hence $S=\emptyset$.
\end{proof}

\begin{remark}
    The argument needs the \emph{least element} $k$ of $S$, not just a
    greatest lower bound: we must know that $k\in S$, i.e.\ that $P(k)$
    fails. For a nonempty subset of $\N$ this is exactly what
    well-ordering provides; for a subset of $\R$ the infimum need not belong
    to the set (\S\ref{sec:bounds}).
\end{remark}

%!TEX root = ../../note.tex
\subsection{The Axioms for the Real Numbers}\label{sec:axioms}

The real numbers form a \term{complete ordered field}. The axioms come in
three groups: the \emph{field axioms} govern $+$ and $\cdot$ (they let us
compute), the \emph{order axioms} govern $<$ (they let us compare), and the
\emph{completeness axiom} is what makes limits possible. The first two
groups hold in $\Q$ as well; only completeness separates $\R$ from $\Q$.

\subsubsection{The field axioms}
\begin{definition}\label{def:field}
    A \term{field} is a set $F$ with two operations
    $+,\,\cdot : F \times F \to F$ such that for all $a,b,c \in F$:
    \begin{enumerate}[label=\textbf{F\arabic*.}, leftmargin=2.5em]
        \item $a + (b + c) = (a + b) + c$;
        \item $a + b = b + a$;
        \item there is $0 \in F$ with $a + 0 = a$;
        \item for each $a$ there is $-a \in F$ with $a + (-a) = 0$;
        \item $a (b c) = (a b) c$;
        \item $a b = b a$;
        \item there is $1 \in F$, $1 \neq 0$, with $a \cdot 1 = a$;
        \item for each $a \neq 0$ there is $a^{-1} \in F$ with $a a^{-1} = 1$;
        \item $a (b + c) = a b + a c$.
    \end{enumerate}
\end{definition}

$\Q$ is a field, so these axioms say nothing about size or order. What they
do is justify every routine algebraic manipulation. Even ``obvious'' facts
must be derived from them:

\begin{theorem}
    Let $a,b,u,z \in \R$.
    \begin{enumerate}[label=(\alph*)]
        \item If $z + a = a$, then $z = 0$.
        \item If $b \neq 0$ and $u b = b$, then $u = 1$.
        \item $a \cdot 0 = 0$.
    \end{enumerate}
\end{theorem}

\begin{proof}
    (a) $z = z + 0 = z + (a - a) = (z + a) - a = a - a = 0$.

    (b) $u = u \cdot 1 = u (b b^{-1}) = (u b) b^{-1} = b b^{-1} = 1$.

    (c) $a\cdot0=a\cdot(0+0)=a\cdot0+a\cdot0$; adding $-(a\cdot0)$ to both
    sides gives $0=a\cdot0$.
\end{proof}

\subsubsection{The order axioms}
\begin{definition}\label{def:order}
    There is a nonempty set $P\subseteq\R$, the \term{positive real
    numbers}, such that
    \begin{enumerate}[label=\textbf{O\arabic*.}, leftmargin=2.5em]
        \item (\emph{trichotomy}) for every $a \in \R$ exactly one of
        $a = 0$, $a \in P$, $-a \in P$ holds;
        \item $a, b \in P \Rightarrow a + b \in P$;
        \item $a, b \in P \Rightarrow a b \in P$.
    \end{enumerate}
    We write $a>b$ (or $b<a$) if $a-b\in P$, and $a\geq b$ (or $b\leq a$)
    if $a-b\in P\cup\set{0}$.
\end{definition}

Every inequality is thus a statement about membership in $P$, and every
manipulation of inequalities reduces to O1--O3. Trichotomy is what makes
case analysis possible, as in the following example.

\begin{example}
    If $a\neq0$, then $a^2>0$.
\end{example}

\begin{proof}
    By trichotomy, $a\in P$ or $-a\in P$. If $a\in P$, then $a^2=a\cdot a\in
    P$ by O3. If $-a\in P$, then $(-a)(-a)\in P$ by O3, and
    $(-a)(-a)=a^2$: multiplying $a+(-a)=0$ by $-a$ gives
    $(-a)a+(-a)(-a)=0$, so $(-a)(-a)=-\bigl((-a)a\bigr)=-(-a^2)=a^2$.
\end{proof}

The positive rationals also satisfy O1--O3, so $\Q$ is an ordered field
too. What it lacks is completeness (\S\ref{sec:completeness}).

\subsubsection{Absolute value}
\begin{definition}
    The \term{absolute value} of $a\in\R$ is
    $\abs{a} = a$ if $a \geq 0$ and $\abs{a} = -a$ if $a < 0$.
\end{definition}

Geometrically, $\abs{a-b}$ is the distance between $a$ and $b$ on the line.

\begin{theorem}\label{thm:abs}
    For all $a, b \in \R$ and $c>0$:
    \begin{enumerate}[label=(\alph*)]
        \item $\abs{ab} = \abs{a}\abs{b}$;
        \item $\abs{a}^2 = a^2$;
        \item $\abs{a} \leq c \iff -c \leq a \leq c$;
        \item $-\abs{a} \leq a \leq \abs{a}$.
    \end{enumerate}
\end{theorem}

\begin{proof}
    Each part follows by splitting into the cases $a\geq0$ and $a<0$.
    Parts (b) and (d) are immediate from the definition.

    (a) If $a,b\geq0$, then $ab\geq0$ and $\abs{ab}=ab=\abs a\abs b$. If
    $a\geq0>b$, then $ab\leq0$ and $\abs{ab}=-ab=a(-b)=\abs a\abs b$;
    the case $b\geq0>a$ is symmetric. If $a,b<0$, then $ab>0$ and
    $\abs{ab}=ab=(-a)(-b)=\abs a\abs b$.

    (c) Let $\abs a\leq c$. If $a\geq0$, then $a=\abs a\leq c$ and
    $-c<0\leq a$. If $a<0$, then $-a=\abs a\leq c$, i.e.\ $-c\leq a$, and
    $a<0<c$. Conversely, let $-c\leq a\leq c$. If $a\geq0$, then
    $\abs a=a\leq c$; if $a<0$, then $-c\leq a$ gives $\abs a=-a\leq c$.
\end{proof}

Part (c) converts one inequality for $\abs{x}$ into two ordinary
inequalities for $x$ and back; it will be used constantly. Its first
application:

\begin{proposition}[Triangle inequality]\label{prop:triangle}
    $\abs{a + b } \leq \abs{a} + \abs{b}$ for all $a, b \in \R$.
\end{proposition}

\begin{proof}
    By Theorem~\ref{thm:abs}(d), $-\abs{a} \leq a \leq \abs{a}$ and
    $-\abs{b} \leq b \leq \abs{b}$. Adding,
    \begin{equation*}
        -\paren{\abs{a} + \abs{b}} \leq a + b \leq \abs{a} + \abs{b}.
    \end{equation*}
    If $\abs a+\abs b=0$, then $a+b=0$ and there is nothing to prove.
    Otherwise Theorem~\ref{thm:abs}(c), with $c=\abs{a}+\abs{b}$, gives
    $\abs{a+b}\leq\abs a+\abs b$.
\end{proof}

\begin{corollary}\label{cor:triangle}
    For all $a,b\in\R$:
    \begin{enumerate}[label=(\arabic*)]
        \item $\abs{\abs{a} - \abs{b}} \leq \abs{a - b}$;
        \item $\abs{a - b} \leq \abs{a} + \abs{b}$.
    \end{enumerate}
\end{corollary}

\begin{proof}
    (1) Apply the triangle inequality to $a=(a-b)+b$ and to $b=(b-a)+a$:
    \begin{equation*}
        \abs a-\abs b\leq\abs{a-b},
        \qquad
        \abs b-\abs a\leq\abs{b-a}=\abs{a-b}.
    \end{equation*}
    Hence $-\abs{a-b}\leq\abs a-\abs b\leq\abs{a-b}$, which is (1) by
    Theorem~\ref{thm:abs}(c).

    (2) $\abs{a - b} = \abs{a + (-b)} \leq \abs{a} + \abs{-b}
    = \abs{a} + \abs{b}$.
\end{proof}

\subsubsection{The completeness axiom}
We state the third axiom now. The notions it uses are made precise in \S\ref{sec:bounds}, and its consequences are the
subject of \S\ref{sec:completeness}.

\begin{axiom}[Completeness]
    Every nonempty subset of $\R$ that is bounded above has a least upper
    bound (supremum) in $\R$.
\end{axiom}

$\Q$ does not have this property: the set $\set{q\in\Q : q^2<2}$ is bounded
above in $\Q$ but has no least upper bound in $\Q$ (see
Proposition~\ref{prop:Q-incomplete}).

%!TEX root = ../../note.tex
\subsection{Bounded Sets; Supremum and Infimum}\label{sec:bounds}

Bounds make sense in any ordered field. What is special about $\R$ is that
a set bounded above always has a \emph{least} upper bound in $\R$; this
section introduces the language needed to say so.

\begin{definition}\label{def:bounded}
    A nonempty set $S \subseteq \R$ is
    \begin{enumerate}[label=(\arabic*)]
        \item \term{bounded above} if there is $u \in \R$ (an \emph{upper
        bound}) with $x \leq u$ for all $x \in S$;
        \item \term{bounded below} if there is $w \in \R$ (a \emph{lower
        bound}) with $w \leq x$ for all $x \in S$;
        \item \term{bounded} if it is bounded above and below, and
        \emph{unbounded} otherwise.
    \end{enumerate}
\end{definition}

\begin{definition}\label{def:sup}
    An upper bound $u_0$ of $S$ is called the \term{supremum} (least upper
    bound) of $S$, written $u_0 = \sup S$, if $u_0 \leq u$ for every upper
    bound $u$ of $S$.

    A lower bound $w_0$ of $S$ is called the \term{infimum} (greatest lower
    bound) of $S$, written $w_0 = \inf S$, if $w_0 \geq w$ for every lower
    bound $w$ of $S$.
\end{definition}

The notation $\sup S$ is justified because a set has at most one supremum:
if $u_1,u_2$ are both suprema, then $u_1\leq u_2$ (as $u_2$ is an upper
bound and $u_1$ is least) and likewise $u_2\leq u_1$. The same holds for
the infimum.

Checking that $u$ is the \emph{least} upper bound requires comparing $u$ with
every other upper bound. The following lemma replaces this by a condition on
$S$ itself: no matter how small a margin $\epsilon$ we allow, $S$ has a point
within $\epsilon$ of $u$.

\begin{lemma}[$\epsilon$-characterization of the supremum]\label{lem:eps-sup}
    An upper bound $u$ of a nonempty set $S\subseteq\R$ is the supremum of
    $S$ if and only if
    \begin{equation*}
        (\forall \epsilon > 0)(\exists x_\epsilon \in S)\,[u - \epsilon < x_\epsilon].
    \end{equation*}
\end{lemma}

\begin{proof}
    ($\Rightarrow$) Let $u=\sup S$ and $\epsilon>0$. Since
    $u-\epsilon<u$ and $u$ is the least upper bound, $u-\epsilon$ is not an
    upper bound, so some $x_\epsilon\in S$ satisfies $u-\epsilon<x_\epsilon$.

    ($\Leftarrow$) If $u$ were not least, there would be an upper bound
    $w<u$. For $\epsilon=u-w>0$ we get $x_\epsilon\in S$ with
    $w=u-\epsilon<x_\epsilon$, so $w$ is not an upper bound after all.
\end{proof}

Thus, to prove $u=\sup S$ one checks two things:
(1) $u$ is an upper bound of $S$;
(2) for every $\epsilon>0$ some $x\in S$ satisfies $u-\epsilon<x\leq u$.
For $\inf S$ reverse the inequalities.

\begin{example}\label{ex:half-open}
    Let $S = \set{x \in \R : 0 \leq x < 1}$. Then $\sup S = 1$ and
    $\inf S = 0$.
\end{example}

\begin{proof}
    \emph{$\sup S=1$.} Clearly $1$ is an upper bound. By completeness
    $u=\sup S$ exists, and $u\leq1$. Suppose $u<1$ and put
    $\epsilon=\frac{1-u}{2}>0$. Then
    \begin{equation*}
        u<u+\epsilon=\frac{u+1}{2}<1,
    \end{equation*}
    and $u+\epsilon\geq0$ because $u\geq0$ (as $0\in S$). So
    $u+\epsilon\in S$ exceeds the upper bound $u$, a contradiction. Hence
    $u=1$.

    \emph{$\inf S=0$.} Clearly $0$ is a lower bound. Every lower bound $w$
    satisfies $w\leq0$ because $0\in S$. Hence $0$ is the greatest lower
    bound.
\end{proof}

\begin{remark}
    In this example $\sup S=1\notin S$, so $S$ has no largest element,
    whereas $\inf S=0\in S$ is also its least element. A supremum is a
    \emph{maximum} exactly when it belongs to the set. Note also that
    ``bounded'' is a property of the set $S$, while $\sup S$ and $\inf S$
    are numbers attached to it.
\end{remark}

%!TEX root = ../../note.tex
\subsection{Consequences of the Completeness Axiom}\label{sec:completeness}

Every result of this section is proved by the same device: form the
supremum of a well-chosen set (the completeness axiom guarantees it exists)
and play the two defining properties of the supremum against each other.

\subsubsection{Infima}
\begin{corollary}\label{cor:inf-exists}
    Every nonempty subset of $\R$ that is bounded below has an infimum
    in $\R$.
\end{corollary}

The idea is to reflect the set through $0$: lower bounds of $S$ become
upper bounds of $-S$.

\begin{proof}
    Let $w$ be a lower bound of $S$ and put $S' = \set{-s : s \in S}$.
    Since $w\leq s$ implies $-s\leq -w$, the number $-w$ is an upper bound of
    $S'$. By completeness $u = \sup S'$ exists. We claim $-u = \inf S$.

    \emph{$-u$ is a lower bound:} $-s\leq u$ for all $s\in S$, so
    $-u\leq s$.

    \emph{$-u$ is the greatest one:} if $w$ is any lower bound of $S$, then
    $-w$ is an upper bound of $S'$, so $u\leq -w$, i.e.\ $w\leq -u$.
\end{proof}

\begin{example}
    The set $S = \set{\paren{1 + \frac{1}{n}}^n : n \in \N}$ satisfies
    $1<x<3$ for all $x\in S$; by completeness, it has a supremum.
\end{example}

\begin{proof}
    By the binomial theorem,
    \begin{equation*}
        \paren{1+\frac1n}^n
        =1+\binom n1\frac1n+\sum_{k=2}^n\binom nk\frac1{n^k}>1.
    \end{equation*}
    For $0\leq k\leq n$,
    $\binom nk=\frac{n(n-1)\cdots(n-k+1)}{k!}\leq\frac{n^k}{k!}$, and
    $k!\geq2^{k-1}$ for $k\geq2$ (each factor $2,\ldots,k$ is at least $2$).
    Therefore
    \begin{equation*}
        \paren{1+\frac1n}^n
        \leq\sum_{k=0}^n\frac1{k!}
        < 1+1+\sum_{k=2}^{\infty}\frac1{2^{k-1}}=3.
    \end{equation*}
\end{proof}

\subsubsection{The Archimedean property}
\begin{theorem}[Archimedean property]\label{thm:archimedes}
    For every $x \in \R$ there is $n_x \in \N$ with $x \leq n_x$.
    In other words, $\N$ is not bounded above in $\R$.
\end{theorem}

The proof uses two ingredients: completeness, to form $u=\sup\N$, and the
closure of $\N$ under $n\mapsto n+1$, to step beyond $u$.

\begin{proof}
    Suppose some $x \in \R$ satisfies $x > n$ for all $n \in \N$. Then $\N$
    is nonempty and bounded above, so by completeness $u = \sup \N$ exists in
    $\R$. Since $u - 1 < u$ and $u$ is the \emph{least} upper bound, $u-1$ is
    not an upper bound of $\N$ (Lemma~\ref{lem:eps-sup} with $\epsilon=1$):
    there is $m \in \N$ with $u - 1 < m$. Then $u < m+1$. But $m + 1 \in \N$
    and $u$ is an upper bound of $\N$, so $m+1\leq u$ --- a contradiction.
\end{proof}

\begin{corollary}\label{cor:archimedes}
    \leavevmode
    \begin{enumerate}[label=(\alph*)]
        \item If $S = \set{1/n : n \in \N}$, then $\inf S = 0$.
        \item If $t > 0$, there is $n_t \in \N$ with $0 < 1/n_t < t$.
        \item If $y > 0$, there is $n_y \in \N$ with $n_y - 1 \leq y < n_y$.
    \end{enumerate}
\end{corollary}

\begin{proof}
    We prove (b) first, since (a) uses it.

    (b) Choose $m\in\N$ with $1/t\leq m$ and put $n_t=m+1$. Then
    $n_t>1/t$, i.e.\ $0<1/n_t<t$.

    (a) $0$ is a lower bound of $S$. If $b>0$, part (b) gives $n$ with
    $1/n<b$, so $b$ is not a lower bound. Hence $\inf S=0$.

    (c) We want the \emph{first} natural number to the right of $y$, so we
    consider all candidates,
    \begin{equation*}
        A=\set{n\in\N:y<n}.
    \end{equation*}
    By the Archimedean property $A\neq\emptyset$, so by well-ordering
    (Theorem~\ref{thm:wop}) $A$ has a least element $n_y$, and $y<n_y$
    because $n_y\in A$. If we had $n_y-1>y$, then $n_y-1>0$, so
    $n_y-1\in\N$ and hence $n_y-1\in A$, contradicting the minimality of
    $n_y$. Thus $n_y-1\leq y<n_y$.
\end{proof}

\begin{remark}
    Part (c) is a \emph{minimum} argument, not a supremum argument. The set
    $A$ is unbounded above, so $\sup A$ does not exist; what we need is the
    least element of $A$, and well-ordering supplies it for any nonempty
    subset of $\N$, bounded or not.
\end{remark}

\subsubsection{Existence of \texorpdfstring{$\sqrt2$}{sqrt 2}}
\begin{proposition}\label{prop:sqrt2-exists}
    There is a real number $x$ with $x^2 = 2$.
\end{proposition}

We define $x$ as the supremum of the set $S$ of all nonnegative $t$ with
$t^2<2$, i.e.\ of all ``approximations of $\sqrt2$ from below''. Then we
rule out both alternatives to $x^2=2$:
\begin{itemize}
    \item if $x^2<2$, then $x$ is too small: $x+\frac1n$ is still in $S$,
    so $x$ is not an upper bound;
    \item if $x^2>2$, then $x$ is too large: $x-\frac1n$ is still an upper
    bound, so $x$ is not the least one.
\end{itemize}
In both cases $n$ is chosen by the Archimedean property.

\begin{proof}
    Let $S = \set{t \in \R : t \geq 0,\ t^2 < 2}$. Then $1 \in S$, and $2$
    is an upper bound of $S$ (if $t>2$ then $t^2>4$). By completeness
    $x = \sup S$ exists, and $x \geq 1$.

    \emph{Case $x^2 < 2$.} For every $n\in\N$, since $1/n^2\leq 1/n$,
    \begin{equation*}
        \paren{x + \frac{1}{n}}^2 = x^2 + \frac{2x}{n} + \frac{1}{n^2}
        \leq x^2 + \frac{2x+1}{n}.
    \end{equation*}
    As $2-x^2>0$, the Archimedean property gives $n$ with
    $\frac{2x+1}{n} < 2 - x^2$. Then $\paren{x + \frac{1}{n}}^2 < 2$, so
    $x+\frac1n\in S$ although $x+\frac1n>x$. This contradicts that $x$ is an
    upper bound of $S$.

    \emph{Case $x^2 > 2$.} For every $n\in\N$,
    \begin{equation*}
        \paren{x - \frac{1}{n}}^2 = x^2 - \frac{2x}{n} + \frac{1}{n^2}
        > x^2 - \frac{2x}{n}.
    \end{equation*}
    As $x^2-2>0$, the Archimedean property gives $n$ with
    $\frac{2x}{n} < x^2 - 2$. Then $\paren{x - \frac{1}{n}}^2 > 2 > t^2$
    for every $t\in S$. Since $x-\frac1n\geq0$ and $s\mapsto s^2$ is strictly
    increasing on $[0,\infty)$, this gives $t<x-\frac1n$ for all $t\in S$.
    So $x-\frac1n$ is an upper bound of $S$ smaller than $x$, contradicting
    $x=\sup S$.

    Hence $x^2 = 2$.
\end{proof}

\subsubsection{Density of \texorpdfstring{$\Q$}{Q} in \texorpdfstring{$\R$}{R}}
\begin{theorem}\label{thm:Q-dense}
    If $x,y\in\R$ and $x < y$, there is $r \in \Q$ with $x < r < y$.
\end{theorem}

Let $0<x<y$ first. We look for $r=m/n$. Choose $n$ so large that the grid
of multiples of $1/n$ is finer than the gap $y-x$; then some grid point must
fall in $(x,y)$, and $m/n$ is the first grid point to the right of $x$.

\begin{proof}
    \emph{Case $0<x<y$.} By the Archimedean property choose $n\in\N$ with
    $1/n<y-x$, i.e.\ $nx+1<ny$. By Corollary~\ref{cor:archimedes}(c) applied
    to $nx>0$, there is $m\in\N$ with $m-1\leq nx<m$. Then
    \begin{equation*}
        nx<m\leq nx+1<ny,
    \end{equation*}
    and dividing by $n$ gives $x<m/n<y$.

    \emph{Case $x\leq0$.} Choose $k\in\N$ with $-x<k$. Then $0<x+k<y+k$, so
    by the first case there is $q\in\Q$ with $x+k<q<y+k$, and $r=q-k\in\Q$
    satisfies $x<r<y$.
\end{proof}

\begin{corollary}
    If $x<y$, there is an irrational $z$ with $x < z < y$.
\end{corollary}

\begin{proof}
    Apply Theorem~\ref{thm:Q-dense} to $x/\sqrt2<y/\sqrt2$ to obtain a
    rational $r\neq0$ between them (if $x<0<y$, take $r$ between $0$ and
    $y/\sqrt2$). Then $z=\sqrt2\,r$ satisfies $x<z<y$, and $z\notin\Q$,
    since otherwise $\sqrt2=z/r$ would be rational.
\end{proof}

\begin{corollary}
    If $x<y$, there are infinitely many rationals between $x$ and $y$.
\end{corollary}

\begin{proof}
    Choose $r_1\in\Q\cap(x,y)$; having chosen $r_k$, choose
    $r_{k+1}\in\Q$ with $r_k<r_{k+1}<y$. The numbers $r_1<r_2<\cdots$ are
    distinct rationals in $(x,y)$.
\end{proof}

We can now make precise the claim that $\Q$ is not complete. The set $T$
below is the rational analogue of the set $S$ used to construct $\sqrt2$;
its supremum ``ought to be'' $\sqrt2$, which is not in $\Q$.

\begin{proposition}\label{prop:Q-incomplete}
    The set $T=\set{r\in\Q: r^2<2}$ is nonempty and bounded above in $\Q$,
    but has no least upper bound in $\Q$.
\end{proposition}

\begin{proof}
    $1\in T$, and $2$ is an upper bound. Let $r\in\Q$ be any upper bound of
    $T$; we find a smaller rational upper bound. Note $r\geq1$ and
    $r\neq\sqrt2$ (Theorem~\ref{thm:sqrt2-irrational}).

    If $r<\sqrt2$, Theorem~\ref{thm:Q-dense} gives $q\in\Q$ with
    $r<q<\sqrt2$; then $q>0$ and $q^2<2$, so $q\in T$ and $q>r$,
    contradicting that $r$ is an upper bound. Hence $r>\sqrt2$, and
    Theorem~\ref{thm:Q-dense} gives $q\in\Q$ with $\sqrt2<q<r$. For $t\in T$
    we have $t^2<2<q^2$ with $q>0$, so $t<q$: thus $q$ is a rational upper
    bound of $T$ smaller than $r$.
\end{proof}

\subsubsection{Countable sets}
Are there more rational or irrational numbers? To give the question a
meaning, we compare sets by means of bijections.

\begin{definition}
    Sets $A$ and $B$ have the same \term{cardinality}, $|A|=|B|$, if there
    is a bijection (a map that is injective and surjective) $f:A\to B$. A set
    $S$ is
    \begin{enumerate}[label=(\arabic*)]
        \item \term{finite} if $S=\emptyset$ or there is a bijection
        $\set{1,\ldots,n}\to S$ for some $n\in\N$;
        \item \term{countably infinite} if there is a bijection $\N\to S$;
        \item \term{countable} if it is finite or countably infinite, and
        \term{uncountable} otherwise.
    \end{enumerate}
\end{definition}

\begin{example}
    (a) $f(n)=2n$ is a bijection from $\N$ onto the even natural numbers.

    (b) $g(1)=0$, $g(2k)=k$, $g(2k+1)=-k$ $(k\in\N)$ is a bijection
    $\N\to\Z$. Thus both sets are countably infinite: a set can have the
    same cardinality as a proper subset of itself.
\end{example}


%!TEX root = ../note.tex
\section{Sequences and Limits}

The completeness of $\R$ is useful only once we can speak of numbers being
approached. The simplest setting for this is a sequence, and the central
notion is that of its limit.

%!TEX root = ../../note.tex
\subsection{The Limit of a Sequence}

\subsubsection{Definition and first examples}
A \term{sequence} of real numbers is a function $X:\N\to\R$. We write
$x_n=X(n)$ for its $n$th term and denote the sequence by
$\set{x_n}_{n=1}^{\infty}$ or simply $\set{x_n}$.

\begin{definition}\label{def:limit}
    A sequence $\set{x_n}$ \term{converges} to $a\in\R$, written
    $\lim_{n\to\infty}x_n=a$ or $x_n\to a$, if
    \begin{equation*}
        (\forall \epsilon>0)(\exists n_\epsilon\in\N)(\forall n\geq n_\epsilon)
        \,\bigl[\abs{x_n-a}<\epsilon\bigr].
    \end{equation*}
\end{definition}

In words: for every tolerance $\epsilon>0$, all sufficiently late terms lie
within $\epsilon$ of $a$. The index $n_\epsilon$ is allowed to depend on
$\epsilon$ (smaller tolerances usually require later indices), and finitely
many early terms may lie anywhere. In the figure, $x_n=1/n$, $a=0$,
$\epsilon=0.3$: every term from $n_\epsilon=4$ on lies in the shaded band.

\begin{center}
\begin{tikzpicture}[x=0.55cm,y=2.4cm,>=Stealth]
    \fill[red!8] (0,0) rectangle (11,0.3);
    \draw[densely dashed,red!75!black] (0,0.3) -- (11,0.3) node[right] {$a+\epsilon$};
    \draw[->] (0,0) -- (11.5,0) node[right] {$n$};
    \draw[->] (0,0) -- (0,1.15) node[above] {$x_n$};
    \node[left] at (0,0) {$a=0$};
    \draw[densely dashed,blue!70!black] (4,0) -- (4,0.95);
    \node[below] at (4,0) {$n_\epsilon=4$};
    \foreach \n in {1,...,10} {
        \fill[teal!70!black] (\n,{1/\n}) circle (1.7pt);
    }
    \node[teal!70!black,anchor=west] at (1.3,0.85) {$x_n=1/n$};
\end{tikzpicture}
\end{center}

The interval
\[
    V_\epsilon(a)=\set{x\in\R:\abs{x-a}<\epsilon}=(a-\epsilon,a+\epsilon)
\]
is called the \term{$\epsilon$-neighbourhood} of $a$. Thus $x_n\to a$ if and
only if every $\epsilon$-neighbourhood of $a$ contains all terms of the
sequence from some index on.

To prove $x_n\to a$ from the definition, one fixes $\epsilon>0$, estimates
$\abs{x_n-a}$ from above by something simple, and chooses $n_\epsilon$ ---
usually by the Archimedean property --- to make that bound less than
$\epsilon$.

\begin{example}\label{ex:1/n}
    $\lim_{n\to\infty}\frac1n=0$.
\end{example}

\begin{proof}
    Given $\epsilon>0$, choose $n_\epsilon$ with $1/n_\epsilon<\epsilon$
    (Corollary~\ref{cor:archimedes}(b)). For $n\geq n_\epsilon$,
    $\abs{\frac1n-0}=\frac1n\leq\frac1{n_\epsilon}<\epsilon$.
\end{proof}

\begin{example}
    $\lim_{n\to\infty}\frac{4n+3}{n+1}=4$.
\end{example}

\begin{proof}
    Given $\epsilon>0$, choose $n_\epsilon$ with $1/n_\epsilon<\epsilon$. For
    $n\geq n_\epsilon$,
    \begin{equation*}
        \abs{\frac{4n+3}{n+1}-4}
        =\frac1{n+1}<\frac1n\leq\frac1{n_\epsilon}<\epsilon.
    \end{equation*}
\end{proof}

\subsubsection{Uniqueness of the limit}
\begin{proposition}\label{prop:eps-equal}
    If $x,y\in\R$ and $\abs{x-y}<\epsilon$ for every $\epsilon>0$, then $x=y$.
\end{proposition}

\begin{proof}
    If $x\ne y$, take $\epsilon=\abs{x-y}/2>0$; then $\abs{x-y}>\epsilon$.
\end{proof}

\begin{theorem}\label{thm:limit-unique}
    A sequence has at most one limit.
\end{theorem}

If $x_n\to a$ and $x_n\to b$, a late enough term is close to both $a$ and
$b$, so by the triangle inequality $a$ and $b$ are close to each other.

\begin{proof}
    Let $\epsilon>0$. Choose $n_1$ with $\abs{x_n-a}<\epsilon/2$ for
    $n\geq n_1$ and $n_2$ with $\abs{x_n-b}<\epsilon/2$ for $n\geq n_2$. For
    $n\geq\max\set{n_1,n_2}$,
    \begin{equation*}
        \abs{a-b}\leq\abs{a-x_n}+\abs{x_n-b}<\frac\epsilon2+\frac\epsilon2=\epsilon.
    \end{equation*}
    Since $\epsilon>0$ was arbitrary, $a=b$ by
    Proposition~\ref{prop:eps-equal}.
\end{proof}

\subsubsection{Divergence and boundedness}
A sequence that does not converge to any real number is said to
\term{diverge}.

\begin{example}\label{ex:0101}
    The sequence $0,1,0,1,\ldots$ ($x_n=0$ for odd $n$, $x_n=1$ for even $n$)
    diverges: both values recur arbitrarily late, so the terms cannot all
    stay near one number.
\end{example}

\begin{proof}
    Suppose $x_n\to a$ and take $\epsilon=1/2$. Choosing an odd and an even
    index $n\geq n_\epsilon$ gives $\abs{0-a}<\frac12$ and $\abs{1-a}<\frac12$,
    i.e.\ $-\frac12<a<\frac12$ and $\frac12<a<\frac32$, which is impossible.
\end{proof}

\begin{definition}
    A sequence $\set{x_n}$ is \term{bounded} if there is $M\in\R$ with
    $\abs{x_n}\leq M$ for all $n\in\N$.
\end{definition}

\begin{theorem}\label{thm:conv-bounded}
    Every convergent sequence is bounded.
\end{theorem}

Convergence controls all terms from some index on; only finitely many terms
remain, and a finite set is bounded.

\begin{proof}
    Let $x_n\to a$. With $\epsilon=1$, choose $n_1$ such that
    $\abs{x_n-a}<1$ for $n\geq n_1$; then
    $\abs{x_n}\leq\abs{x_n-a}+\abs a<1+\abs a$ for these $n$. Hence
    \begin{equation*}
        M=\max\set{\abs{x_1},\ldots,\abs{x_{n_1-1}},\,1+\abs{a}}
    \end{equation*}
    bounds every term (if $n_1=1$, take $M=1+\abs a$).
\end{proof}

The converse is false: the sequence of Example~\ref{ex:0101} is bounded but
divergent. On the other hand, by Theorem~\ref{thm:conv-bounded} an
\emph{unbounded} sequence must diverge.

\begin{example}
    The following sequences are unbounded, hence divergent.

    (a) $a_n=\frac{n^2+1}{n}$. Here $\abs{a_n}=n+\frac1n$. Given $M>0$,
    choose $n_M\in\N$ with $n_M>M$; then $\abs{a_{n_M}}>M$.

    (b) $a_n=(-1)^n n$. Here $\abs{a_n}=n$, and again $\abs{a_{n_M}}>M$
    whenever $n_M>M$.
\end{example}

\subsubsection{Infinite limits}
\begin{definition}
    A sequence $\set{x_n}$ \term{diverges to $+\infty$}, written
    $\lim_{n\to\infty}x_n=+\infty$, if for every $M>0$ there is
    $n_M\in\N$ such that $x_n>M$ for all $n\geq n_M$. It diverges to
    $-\infty$ if for every $M>0$ there is $n_M$ such that $x_n<-M$ for all
    $n\geq n_M$.
\end{definition}

Such a sequence is unbounded and therefore \emph{divergent}; the notation
$\lim x_n=\pm\infty$ only records the manner in which it diverges.

\begin{example}
    $\lim_{n\to\infty}\frac{n^2}{1-n}=-\infty$.
\end{example}

\begin{proof}
    For $n\geq2$ we have $0<n-1<n$, so
    $\frac{n^2}{n-1}>\frac{n^2}{n}=n$, i.e.\ $\frac{n^2}{1-n}<-n$. Given
    $M>0$, choose $n_M\in\N$ with $n_M\geq2$ and $n_M>M$. For $n\geq n_M$,
    \begin{equation*}
        \frac{n^2}{1-n}<-n\leq-n_M<-M.
    \end{equation*}
\end{proof}

\begin{example}
    $\lim_{n\to\infty}\paren{\sqrt{n}+1}=+\infty$.
\end{example}

\begin{proof}
    Given $M>0$, we must make $\sqrt n>M-1$ for all late $n$. Choose
    $n_M\in\N$ with $n_M>M^2$. For $n\geq n_M$,
    $\sqrt n\geq\sqrt{n_M}>M>M-1$, so $\sqrt n+1>M$.
\end{proof}

\subsubsection{Limits and arithmetic operations}
% The lecture stops after the first two parts; the theorem is continued
% (sums, products, quotients of limits) in the next lecture.
\begin{theorem}\label{thm:limit-laws}
    Let $\lim_{n\to\infty} x_n = a$ and $c\in\R$. Then
    \begin{enumerate}[label=(\arabic*)]
        \item $\lim_{n \to \infty} c = c$ (the constant sequence);
        \item $\lim_{n \to \infty} cx_n = ca$.
    \end{enumerate}
\end{theorem}

\begin{proof}
    (1) $\abs{c-c}=0<\epsilon$ for every $n$.

    (2) If $c=0$, this is (1). If $c\neq0$, given $\epsilon>0$ choose
    $n_\epsilon$ with $\abs{x_n-a}<\epsilon/\abs c$ for $n\geq n_\epsilon$;
    then $\abs{cx_n-ca}=\abs c\abs{x_n-a}<\epsilon$.
\end{proof}



\printindex

\end{document}
